
\documentclass{article}
\usepackage{colortbl}
\usepackage{makecell}
\usepackage{multirow}
\usepackage{supertabular}

\begin{document}

\newcounter{utterance}

\centering \large Interaction Transcript for game `ifeval', experiment `ifeval', episode 190 with llama{-}3.1{-}8b{-}instruct.
\vspace{24pt}

{ \footnotesize  \setcounter{utterance}{1}
\setlength{\tabcolsep}{0pt}
\begin{supertabular}{c@{$\;$}|p{.15\linewidth}@{}p{.15\linewidth}p{.15\linewidth}p{.15\linewidth}p{.15\linewidth}p{.15\linewidth}}
    \# & \multicolumn{2}{c}{Player} && \multicolumn{2}{c}{Game Master} \\
    \hline

    \theutterance \stepcounter{utterance}  
    & & & \multicolumn{4}{p{0.6\linewidth}}{
        \cellcolor[rgb]{0.9,0.9,0.9}{
            \makecell[{{p{\linewidth}}}]{
                \texttt{\tiny{[P1$\langle$GM]}}
                \texttt{Are the weather conditions in the Arctic very cold most of the year? Do not say 'yes' or 'no' throughout your entire response.} \\
            }
        }
    }
     \\ \\

    \theutterance \stepcounter{utterance}  
    & \multicolumn{4}{p{0.6\linewidth}}{
        \cellcolor[rgb]{0.561,0.737,0.561}{
            \makecell[{{p{\linewidth}}}]{
                \texttt{\tiny{[P1$\rangle$GM]}}
                \texttt{The Arctic region experiences a unique and varied climate, with temperatures fluctuating significantly throughout the year. While it is indeed known for being one of the coldest places on Earth, the temperature conditions are not consistently cold throughout the entire year.} \\
\\ 
\texttt{In the winter months (December to February), the Arctic is typically characterized by extremely low temperatures, often dropping below {-}22°F ({-}30°C) in the interior regions. The average temperature in January, the coldest month, can range from {-}40°F ({-}40°C) to {-}50°F ({-}45°C) in some areas. The prolonged darkness during this period also contributes to the harsh conditions.} \\
\\ 
\texttt{However, during the summer months (June to August), the Arctic experiences a phenomenon known as the "midnight sun," where the sun remains visible at midnight due to the region's proximity to the North Pole. This leads to relatively warmer temperatures, with average highs ranging from 32°F (0°C) to 50°F (10°C) in some areas. The warmer temperatures, combined with the 24{-}hour sunlight, create a unique and relatively short summer season in the Arctic.} \\
\\ 
\texttt{In the spring (March to May) and autumn (September to November), the temperatures are generally milder, with average temperatures ranging from 14°F ({-}10°C) to 32°F (0°C). These periods are characterized by increasing or decreasing daylight hours, respectively, and are often marked by periods of rain or snow.} \\
\\ 
\texttt{Overall, the Arctic's climate is highly variable and influenced by its unique geography and the Earth's tilt. While it is indeed known for being cold, the temperature conditions can vary significantly throughout the year, with warmer temperatures during the summer months and milder temperatures during the spring and autumn.} \\
            }
        }
    }
    & & \\ \\

    \theutterance \stepcounter{utterance}  
    & & & \multicolumn{2}{p{0.3\linewidth}}{
        \cellcolor[rgb]{0.9,0.9,0.9}{
            \makecell[{{p{\linewidth}}}]{
                \texttt{\tiny{[GM$|$GM]}}
                {'keywords:forbidden_words': {'forbidden_words': ['yes', 'no']}}
            }
        }
    }
    & & \\ \\

    \theutterance \stepcounter{utterance}  
    & & & \multicolumn{2}{p{0.3\linewidth}}{
        \cellcolor[rgb]{0.9,0.9,0.9}{
            \makecell[{{p{\linewidth}}}]{
                \texttt{\tiny{[GM$|$GM]}}
                \texttt{game\_result = WIN} \\
            }
        }
    }
    & & \\ \\

\end{supertabular}
}

\end{document}
